\documentclass[spanish,letterpaper,oneside]{article}

% Datos del documento
\def\documenttitle {Sociedades del siglo XXI y transformación de la enseñanza}
\def\documentsubtitle {Fundamentos para la enseñanza y el aprendizaje I}
\def\documentsubject {Actividad 3 · Doce argumentos y conclusión personal}
\def\documentauthor {Mtro. Martín Jonathan de la Cruz Muñoz}
\def\coursename {Fundamentos para la enseñanza y el aprendizaje I}
\def\coursecode {MAEM}

% Identidad institucional
\def\universityname {Instituto de Investigación, Innovación y Estudios de Posgrado para la Educación del Estado de Nuevo León}
\def\universityfaculty {Maestría en Aprendizaje y Enseñanza de las Matemáticas}
\def\universitydepartment {Fundamentos para la enseñanza y el aprendizaje I}
\def\universitydepartmentimage {img/iiiepe/marca-iiiepe.png}
\def\universitydepartmentimagecfg {height=1.75cm}
\def\universitylocation {Monterrey, Nuevo León}
\def\studentfullname {Mtro. Martín Jonathan de la Cruz Muñoz}
\def\studentid {00874}
\def\studentemail {martin.delacruz@campus.iiiepe.edu.mx}
\def\programperiod {1er tetramestre, marzo--julio de 2026}
\def\programmodality {Modalidad presencial}
\def\programshort {MAEM}

\def\authortable {
  \begin{tabular}{ll}
    \textbf{Alumno:} & \studentfullname \\
    Matrícula: & \studentid \\
    Correo institucional: & \studentemail \\
    Unidad: & \coursename \\
    Programa: & \programshort \\
    Periodo: & \programperiod \\
    & \\
    \multicolumn{2}{l}{Fecha de realización: 13 de octubre de 2026} \\
    \multicolumn{2}{l}{\universitylocation}
  \end{tabular}
}

\input{template}

\usepackage{helvet}
\renewcommand{\familydefault}{\sfdefault}
\usepackage{array}
\usepackage{longtable}
\usepackage{pdflscape}
\usepackage{tikz}
\usepackage[most]{tcolorbox}
\usetikzlibrary{arrows.meta,positioning,calc,fit,backgrounds}

\definecolor{iiiepeNavy}{HTML}{163B5C}
\definecolor{iiiepeBlue}{HTML}{1F6F9F}
\definecolor{iiiepeAqua}{HTML}{20B7A8}
\definecolor{iiiepeOrange}{HTML}{F28C28}
\definecolor{iiiepePaper}{HTML}{F6F8F7}
\definecolor{iiiepeInk}{HTML}{1E2A32}
\def\portraittitlecolor {iiiepeNavy}
\setcitestyle{authoryear,open={(},close={)}}

\newcommand{\insertcoverwatermark}{%
  \AddToShipoutPictureBG*{%
    \begin{tikzpicture}[remember picture,overlay]
      \node[opacity=.09,inner sep=0pt] at (current page.center)
      {\includegraphics[width=.58\paperwidth]{img/iiiepe/logo-puntos.png}};
    \end{tikzpicture}%
  }%
}

\newtcolorbox{reportelecturabox}[1][]{%
  enhanced,breakable,
  colback=iiiepePaper,colframe=iiiepeNavy,
  boxrule=1pt,leftrule=4pt,arc=2pt,
  left=8pt,right=8pt,top=6pt,bottom=6pt,
  before skip=6pt,after skip=10pt,
  fonttitle=\bfseries\color{white},coltitle=white,
  colbacktitle=iiiepeNavy,
  attach boxed title to top left={xshift=8pt,yshift=-2pt},
  boxed title style={colback=iiiepeNavy,colframe=iiiepeNavy,arc=1pt},
  title={Reporte de lectura},#1}

\newtcolorbox{fichalecturabox}[1][]{%
  enhanced,breakable,
  colback=iiiepeAqua!7,colframe=iiiepeAqua!70!iiiepeNavy,
  boxrule=.6pt,arc=1pt,
  left=7pt,right=7pt,top=4pt,bottom=4pt,
  fontupper=\small,before upper={\setstretch{1.0}},
  fonttitle=\bfseries\color{iiiepeNavy},coltitle=iiiepeNavy,
  colbacktitle=iiiepeAqua!7,
  title={Ficha de la lectura},#1}



\usepackage{tabularx,colortbl,pgfplots,enumitem,needspace}\pgfplotsset{compat=1.18}\definecolor{iiiepePale}{HTML}{EEF5F8}
\setlength{\parskip}{5pt}\setlength{\parindent}{0pt}
\begin{document}\insertcoverwatermark\templatePortrait\templatePagecfg\templateFinalcfg\fancyfoot{}\fancyfoot[L]{\scriptsize IIIEPE · MAEM}\fancyfoot[C]{\scriptsize\thepage}\fancyfoot[R]{\scriptsize Actividad 3}\setstretch{1.3}\raggedbottom
\section{Transformar la enseñanza con criterio}
La lectura de Scott (2015a) permite analizar doce factores que presionan a los sistemas educativos. Mi postura es que la transformación debe cambiar las oportunidades de comprender, participar y argumentar. La incorporación de herramientas digitales o de vocabulario industrial no basta si la actividad intelectual del estudiante sigue limitada a reproducir instrucciones. Los argumentos siguientes relacionan cada factor con mi docencia en formación técnica e ingenierías; son reflexiones propias a partir de la lectura.
\subsection{Competencias nuevas para un mundo complejo}
Los problemas de mantenimiento rara vez se presentan con datos completos y una única ruta. La enseñanza debe incluir selección de información, formulación de supuestos y valoración de consecuencias. Propondría que el estudiante explique por qué un modelo es pertinente antes de resolverlo. La destreza operativa sigue siendo necesaria, pero debe integrarse con juicio y responsabilidad.
\subsection{Nuevas características de las y los estudiantes}
No todos los estudiantes comparten acceso digital, experiencias ni disponibilidad de tiempo. Suponer que una generación domina la tecnología puede ocultar desigualdades de uso y comprensión. Necesito diagnosticar conocimientos y condiciones, ofrecer alternativas y evitar que la familiaridad con una herramienta sustituya la valoración del razonamiento.
\subsection{Desinterés y abandono prematuro de la educación}
Una clase puede perder sentido cuando el estudiante no reconoce qué puede comprender o decidir con lo aprendido. Vincular las matemáticas con un problema de taller ayuda, siempre que la participación sea posible y no una demostración reservada a quienes ya saben. El abandono tiene múltiples causas; no lo reduciría a motivación individual ni atribuiría su solución a una metodología aislada.
\subsection{Valor de la educación formal}
El acceso a información fuera de la escuela cuestiona una enseñanza basada solo en transmitirla. La institución aporta continuidad, diálogo, criterios para evaluar fuentes y oportunidades de participar en conocimiento compartido. En ingeniería, distinguir un tutorial de una explicación justificada es parte de ese valor: la educación formal debe enseñar a comprobar y contextualizar.
\subsection{Escasez de preparación}
La falta de preparación puede aparecer cuando se dominan procedimientos aislados pero no se transfieren a situaciones nuevas. La respuesta docente consiste en conectar representaciones, explicitar supuestos y revisar obstáculos. También implica reconocer necesidades de formación propias: mi experiencia industrial no sustituye conocimiento didáctico sobre cómo se construye una noción matemática.
\subsection{Evolución del mercado de trabajo}
La automatización modifica tareas, pero la educación no debe reducir a la persona a un puesto laboral. En mantenimiento y mecatrónica, aprender a interpretar datos y argumentar una intervención permite actuar ante cambios técnicos; también ayuda a deliberar sobre seguridad, impacto y condiciones de trabajo. Enseñaría competencias transferibles sin prometer empleos ni anticipar tendencias específicas no verificadas.
\subsection{Cambios en modelos de enseñanza y aprendizaje}
Un proyecto modifica el aprendizaje solo si cambia lo que hacen y comprenden los estudiantes. Debe incluir problemas, apoyos, explicación y evaluación, no únicamente una actividad manual. Mi transición hacia una enseñanza humanista exige revisar quién toma decisiones y cómo se usa el error, manteniendo criterios de precisión y comprensión.
\subsection{Nuevos medios didácticos y otras vías de aprendizaje}
Un simulador permite variar parámetros y visualizar relaciones, pero también puede producir respuestas sin comprensión. Lo integraría con predicción, contraste y explicación, junto con una alternativa de papel o material concreto. La elección del medio debe responder a una dificultad de aprendizaje y a las condiciones de acceso, no a su novedad.
\subsection{Evaluación y rendición de cuentas}
Los indicadores ayudan a observar resultados, pero pueden orientar la clase hacia lo que se mide más fácilmente. Evaluaría procedimientos y argumentos, conservaría revisiones y explicaría los criterios. Un promedio no describe todas las condiciones ni demuestra la causa de un desempeño; la rendición de cuentas debe apoyar decisiones y no convertirse en etiquetamiento.
\subsection{Aprendizaje en cualquier momento y lugar}
La continuidad fuera del aula ofrece oportunidades de consulta y práctica, pero no elimina la necesidad de mediación. Una actividad asincrónica requiere instrucciones comprensibles, recursos accesibles y un canal para resolver dudas. También debe respetar que estudiantes de formación técnica o superior pueden tener responsabilidades laborales y familiares.
\subsection{Crecimiento acelerado de la información}
El problema educativo no es únicamente encontrar información, sino discriminarla. Pediría verificar unidades, procedencia, fecha y condiciones de una afirmación técnica. Al comparar dos fórmulas o recomendaciones, el estudiante debe justificar cuál corresponde al problema; acumular enlaces no equivale a construir conocimiento.
\subsection{Del consumidor al productor de conocimiento}
Producir conocimiento en clase puede consistir en formular un modelo, diseñar una prueba y comunicar sus límites. Un estudiante que explica por qué una calibración falla está construyendo una interpretación verificable. La producción debe ser acompañada y revisable; publicar un documento o un prototipo no garantiza por sí solo rigor ni autoría intelectual.
\section{Tensiones entre los factores y decisiones docentes}
La lectura reúne transformaciones sociales y educativas, pero esas tendencias no determinan una única respuesta pedagógica. El acceso permanente a información puede ampliar oportunidades y también dispersar la atención; la preparación para el trabajo puede ofrecer pertinencia y también estrechar la finalidad educativa; la rendición de cuentas puede aportar evidencia y también incentivar una enseñanza limitada al indicador. Mi argumento es que cada factor exige un criterio de decisión, no una adhesión automática a la novedad.
Scott (2015a, pp. 3--9) organiza los factores desarrollados en este análisis. A partir de ese marco, en formación técnica priorizaría tareas que conecten interpretación y comprobación. Por ejemplo, antes de elegir una fórmula de caudal, pediría definir unidades y condiciones; antes de usar una recomendación encontrada en línea, revisar procedencia y ámbito de aplicación; antes de calificar un prototipo, solicitar una explicación individual de una decisión matemática.
Esta postura se diferencia de una adaptación puramente instrumental: la persona debe poder comprender y cuestionar el uso del conocimiento. La relevancia profesional es una vía para plantear problemas, mientras el acompañamiento sostiene participación y juicio. Evaluaría esa relación mediante explicaciones iniciales, versiones corregidas y un caso nuevo, con criterios que se conozcan antes de la tarea.

\section{Conclusión personal}
Mi formación industrial me ha enseñado a organizar y verificar el trabajo; la maestría me lleva a ampliar qué considero una evidencia de aprendizaje. Sostengo que la transformación consiste en conservar rigor y reconocer a la persona que aprende: sus preguntas, condiciones, interpretaciones y posibilidades de participar. Aplicaría esta postura mediante problemas con supuestos explícitos, apoyos diferenciados y revisiones argumentadas. El criterio de mejora será qué logra comprender y transferir cada estudiante, no cuántas herramientas o productos se incorporaron.

\section{Respaldo documental de mi enfoque docente}
El programa de Pensamiento Matemático I de la SEP (2023, pp. 3--4) integra interpretación, argumentación, modelación y comunicación, además del procedimiento. Lo utilizo como referencia documental para mi acercamiento humanista: la precisión técnica adquiere sentido educativo cuando el estudiante puede explicar una decisión y revisar sus razones. Esta relación fundamenta las aplicaciones ilustrativas del reporte; no las convierte en episodios de aula observados.
\needspace{5\baselineskip}\section*{Referencias}
{\setstretch{1.15}\setlength{\parskip}{6pt}
\par\hangindent=1.27cm\hangafter=1  Scott, C. L. (2015a). \textit{El futuro del aprendizaje 1: ¿Por qué deben cambiar el contenido y los métodos de aprendizaje en el siglo XXI?} (Investigación y prospectiva en educación, 13). UNESCO.
\par\hangindent=1.27cm\hangafter=1  Secretaría de Educación Pública. (2022). \textit{Plan de Estudio para la educación preescolar, primaria y secundaria}. SEP.
\par\hangindent=1.27cm Secretaría de Educación Pública. (2023). \textit{Programa de estudios del recurso sociocognitivo Pensamiento Matemático I}. Subsecretaría de Educación Media Superior. \url{https://dgb.sep.gob.mx/storage/recursos/2025/09/nfhyrKiRiF-Pensamiento-Matematico-I.pdf}
}

\end{document}
